
\begin{prop}\label{prop:q-ps}
    For any $n \ge 2$ and $D \subseteq [n-1]$, the specialization of $\Qtil_{n,D}$ given by setting $x_i = q^{i-1}, y_i = tq^{i-1}$ is
    \[
    \Qtil_{n,D}(1,q,q^2,\ldots;t,tq,tq^2,\ldots) = \frac{1}{(q;q)_n} \sum_{S \subseteq [n]} q^{\comaj(D,S)}t^{|S|},
    \]
    where $(a;q)_n = (1-a)(1-aq) \cdots (1-aq^{n-1})$ denotes the $q$-Pochhammer symbol.
\end{prop}
\begin{proof}
    Given a sequence $\mathbf{a} = (a_1 \le a_2 \le \cdots \le a_n)$ from the alphabet $\Acal$ which satisfies:
    \begin{equation}\label{a-con}
a_i = a_{i+1} \in \Acal_+ \Rightarrow i \not\in D, \quad a_i = a_{i+1} \in \Acal_- \Rightarrow i \in D,
\end{equation}
write $b_1 \le b_2 \le \cdots \le b_n$ for the image under projection to $\Acal_+$, and let $S_\mathbf{a} = \{ i : a_i \in \Acal_- \}$. This sequence contributes the term $z_{a_1}z_{a_2} \cdots z_{a_n}$ to $\Qtil_{n,D}(\xb;\yb)$, so upon specializing, this term becomes $q^{b_1 + b_2 + \cdots + b_n - n}t^{|S_\mathbf{a}|}$. 

We claim that \eqref{a-con} is equivalent to the condition that:
    \begin{equation}\label{b-con}
    b_i < b_{i+1} \text{ whenever } i \in D, i+1 \not\in S_\mathbf{a}, \text{ or } i \not\in D, i \in S_\mathbf{a}.
    \end{equation}
    
    Suppose that \eqref{a-con} holds. First assume that $i \in D$ and $i+1 \not\in S_\mathbf{a}$, so that $a_{i+1} \in \Acal_+$. Then we must have either $a_i < a_{i+1}$ or $a_i = a_{i+1} \in \Acal_+$. The latter possibility is not allowed by \eqref{a-con}, so we must have $a_i < a_{i+1}$. Thus $b_i \le a_i < a_{i+1} = b_{i+1}$, so $b_i < b_{i+1}$. Now assume that $i \not\in D$ and $i \in S_\mathbf{a}$. Then $b_i < a_i < a_{i+1}$, so $b_i < b_{i+1}$. Thus \eqref{b-con} holds.
    
    Now assume that \eqref{b-con} holds. Suppose $a_i = a_{i+1} \in \Acal_+$, so that $b_i = b_{i+1}$ and $i+1 \not\in S_\mathbf{a}$. Then by assumption, we must have $i \not\in D$, as required by \eqref{a-con}. On the other hand, if $a_i = a_{i+1} \in \Acal_-$, then again we have $b_i = b_{i+1}$ and $i \in S_\mathbf{a}$, which forces $i \in D$. Therefore \eqref{a-con} and \eqref{b-con} are equivalent.

    Now, for each $1 \le j \le n$, set 
    \[
    r_j = b_j - 1 - |\{ i < j : i \in D, i+1 \not\in S_\mathbf{a} \text{ or } i \not\in D, i \in S_\mathbf{a} \}|.
    \]
    Then $0 \le r_1 \le \cdots \le r_n$ and
    \begin{align*}
    b_1 + \cdots + b_n - n 
    %&= r_1 + \cdots + r_n + \sum_{j=1}^n |\{ i < j : i \in D, i+1 \not\in S_\mathbf{a} \text{ or } i \not\in D, i \in S_\mathbf{a} \}| \\
    &= r_1 + \cdots + r_n + \comaj(D,S_\mathbf{a}),
    \end{align*}
    so that
    \begin{align*}
    \Qtil_{n,D}(1,q,q^2,\ldots;t,tq,tq^2,\ldots) &= \sum_{\substack{\mathbf{a} = (a_1 \le \cdots \le a_n) \\ \text{s.t. } \eqref{a-con} \text{ holds}}} q^{b_1 + b_2 + \cdots + b_n - n}t^{|S_\mathbf{a}|} \\
    &= \sum_{S \subseteq [n]} t^{|S|} \sum_{\substack{\mathbf{a} = (a_1 \le \cdots \le a_n) \\ \text{s.t. } \eqref{a-con} \text{ holds and } S_\mathbf{a} = S}} q^{b_1 + b_2 + \cdots + b_n - n} \\
    &= \sum_{S \subseteq [n]} t^{|S|} q^{\comaj(D,S)} \sum_{0 \le r_1 \le \cdots \le r_n} q^{r_1+\cdots+r_n} \\
    &= \frac{1}{(q;q)_n} \sum_{S \subseteq [n]} q^{\comaj(D,S)} t^{|S|}.\qedhere
    \end{align*}
\end{proof}

The analogous specialization for the super Schur functions $\stil_\lambda(\xb;\yb)$ is given by Macdonald in terms of the \emph{hook lengths} $h(r,c)$ of the cells $(r,c) \in \lambda$. See also \cite{Kratt} for a bijective proof of the following identity.
\begin{thm}[{\cite[p.~27, Ex.~5 and p.~45, Ex.~3]{Mac}}]\label{qt-hook}
   The specialization of $\stil_\lambda(\xb;\yb)$ given by setting $x_i = q^{i-1}$, $y_i = tq^{i-1}$ is given by
   \[
   \stil_\lambda(1,q,q^2,\ldots;t,tq,tq^2,\ldots) = \prod_{(r,c) \in \lambda} \frac{q^{r-1} + tq^{c-1}}{1 - q^{h(r,c)}}.
   \]
\end{thm}
In fact, (a renormalization of) the principal specialization of $\stil_\lambda(\xb;\yb)$ may also be written as the $q,t$-generating function of the statistics $\mathrm{(co)maj}$ and $\negg$:
\begin{thm}\label{thm:s-ps}
For $\lambda \vdash n$, we have
\begin{align*}
 \stil_\lambda(1,q,q^2,\ldots;t,tq,tq^2,\ldots)
   &= \frac{1}{(q;q)_n} \sum_{\mathcal{T} \in \SYT_\pm(\lambda)} q^{\comaj(T)} t^{\negg(T)} \\
   &= \frac{1}{(q;q)_n} \sum_{\mathcal{T} \in \SYT_\pm(\lambda)} q^{\maj(T)} t^{\negg(T)}.
\end{align*}
\end{thm}
\begin{proof}
    We begin by proving the first equality. We have by definition that $\stil_\lambda(\xb;\yb) = \sum_{T \in \SYT(\lambda)} \Qtil_{n,\Des(T)}(\xb;\yb)$, so
    \begin{align*}
        \stil_\lambda(1,q,q^2,\ldots;t,tq,tq^2,\ldots) &= \sum_{T \in \SYT(\lambda)} \Qtil_{n,\Des(T)}(1,q,q^2,\ldots;t,tq,tq^2,\ldots) \\
        &= \frac{1}{(q;q)_n} \sum_{T \in \SYT(\lambda)} \sum_{S \subseteq [n]} q^{\comaj(\Des(T),S)} t^{|S|}
    \end{align*}
    by \Cref{prop:q-ps}. Note that a standard super tableau $\mathcal{T}$ is determined by a standard tableau $T = \mathcal{T}_+$ and a set of negative entries $S = \Neg(\mathcal{T})$. We then have $\comaj(\mathcal{T}) = \comaj(\Des(T),S)$, so the first equality follows.

    Now, for any $D \subseteq [n-1]$ and $S \subseteq [n]$, let 
    \[
    D^* = \{ n-i : i \in D \} \quad \text{and} \quad S^* = \{ n+1-i : i \in S \}
    \]
    denote the \emph{reverse} of $D$ and $S$, respectively. It follows from \cite[Prop.~7.19.2]{Stan} that
    \[
    |\{ T \in \SYT(\lambda) : \Des(T) = D \}| = |\{ T \in \SYT(\lambda) : \Des(T) = D^* \}|.
    \]
    Note furthermore that for any $T \in \SYT(\lambda)$, its transpose $T' \in \SYT(\lambda')$ satisfies $\Des(T') = [n-1] \backslash \Des(T)$. Therefore, there exists a bijection $\varphi : \SYT(\lambda) \to \SYT(\lambda')$ such that
    \[
    \Des(\varphi(T)) = [n-1] \backslash \Des(T)^*
    \]
    for any $T \in \SYT(\lambda)$, which first reverses the descent set of $T$ and then transposes. 
    
    Then $\varphi$ lifts to a bijection $\widetilde{\varphi} : \SYT_\pm(\lambda) \to \SYT_\pm(\lambda')$ which satisfies
    \[
    \Des(\varphi(\mathcal{T})_+) = [n-1] \backslash \Des(\mathcal{T}_+)^* \quad \text{and} \quad \Neg(\varphi(\mathcal{T})) = [n] \backslash \Neg(\mathcal{T})^*
    \]
    for any $\mathcal{T} \in \SYT_\pm(\lambda)$, given by performing $\varphi$ on $\mathcal{T}_+$ and then reverse complementing the set of negative entries. From this, we obtain
    \begin{align*}
    \comaj(\Des(\varphi(\mathcal{T})_+),\Neg(\varphi(\mathcal{T})))
    &= \maj(\Des(\mathcal{T}_+),\Neg(\mathcal{T})) \\
    \negg(\varphi(\mathcal{T}))
    &= n - \negg(\mathcal{T}).
    \end{align*}
    It follows from \Cref{qt-hook} that
    \[
        \stil_{\lambda'}(1,q,q^2,\ldots;t,tq,tq^2,\ldots) = \prod_{(r,c) \in \lambda} \frac{q^{c-1} + tq^{r-1}}{1-q^{h(r,c)}},
    \]
    so that
    \begin{align*}
        \stil_\lambda(1,q,q^2,\ldots;t,tq,tq^2,\ldots) &= t^n\stil_{\lambda'}(1,q,q^2,\ldots;t^{-1},t^{-1}q,t^{-1}q^2,\ldots) \\
        &= t^n \sum_{T' \in \SYT(\lambda')} \Qtil_{n,\Des(T')}(1,q,q^2,\ldots;t^{-1},t^{-1}q,t^{-1}q^2,\ldots) \\
        &= \frac{t^n}{(q;q)_n} \sum_{T' \in \SYT(\lambda')} \sum_{S \subseteq [n]} q^{\comaj(\Des(T'),S)} t^{-|S|} \\
        &= \frac{t^n}{(q;q)_n} \sum_{T \in \SYT(\lambda)} \sum_{S \subseteq [n]} q^{\maj(\Des(T),S)} t^{|S|-n} \\
        &= \frac{1}{(q;q)_n} \sum_{T \in \SYT(\lambda)} \sum_{S \subseteq [n]} q^{\maj(\Des(T),S)} t^{|S|} \\
        &= \frac{1}{(q;q)_n} \sum_{\mathcal{T} \in \SYT_\pm(\lambda)} q^{\maj(\mathcal{T})}t^{\negg(\mathcal{T})}. \qedhere
    \end{align*}
\end{proof}

\begin{rema}
    We are not aware of an explicit description of the map $\varphi$ in the proof of \Cref{thm:s-ps}. See \cite[\S4]{BKS} for a description of the unique $T \in \SYT(\lambda)$ with minimal and maximal major index.
\end{rema}

\Cref{thm:maj-neg-hook} now follows easily from \Cref{thm:s-ps}.
\begin{proof}[Proof (of \Cref{thm:maj-neg-hook}).]
    By \Cref{thm:s-ps},
    \begin{align*}
        \sum_{\mathcal{T} \in \SYT_\pm(\lambda)} q^{\maj(\mathcal{T})}t^{\negg(\mathcal{T})} &= (q;q)_n \stil_\lambda(1,q,q^2,\ldots;t,tq,tq^2,\ldots) \\
        &= (1-q) \cdots (1-q^n) \prod_{(r,c) \in \lambda} \frac{q^{r-1} + tq^{c-1}}{1-q^{h(r,c)}} \\
        &= \frac{1-q}{1-q} \cdots \frac{1-q^n}{1-q} \prod_{(r,c) \in \lambda} \frac{(q^{r-1} + tq^{c-1})(1-q)}{1-q^{h(r,c)}} \\
        &= [n]_q! \prod_{(r,c) \in \lambda} \frac{q^{r-1} + tq^{c-1}}{[h(r,c)]_q}. \qedhere
    \end{align*}
\end{proof}

\section{Irreducible decomposition of \texorpdfstring{$\Ltil_{n,m}$}{Ltilde-n,m}}\label{sec:super-KW}
We now prove \Cref{thm:super-KW}, generalizing Kr{\'a}skiewicz--Weyman's result (\Cref{thm:KW}) in the classical case. Our proof is predominantly algebraic. By manipulating the $q,t$-hook formula (\Cref{thm:maj-neg-hook}), the Schur expansion of $\Ch(\Ltil_{n,m})$ is obtained using the generalized Cauchy identity (\Cref{lem:s-p}), along with the following technical results.

Let $\Lambda_t(\xb)$ denote the ring of symmetric functions in $\xb$ with coefficients in $\CC(t)$. For any $r \ge 1$ define an operator $\Omega^r_1 : \Lambda_t(\xb)[[q]] \to \Lambda_t(\xb)$ by:
\[
\Omega^r_1f(\xb;q,t) = \frac{1}{r} \sum_{\zeta^r = 1} \zeta^{-1} f(\xb;\zeta,t),
\]
summing over all $r$-th roots of unity $\zeta$. Then $\Omega^r_1$ extracts the terms from $f(\xb;q,t)$ whose $q$-exponents are congruent to $1$ modulo $r$; that is,
\[
\Omega^r_1f(\xb;q,t) = \sum_{\ell \in \ZZ} [q^{1+\ell r}]f(\xb;q,t).
\]
We will also require the following lemma, which can be found, for instance, in \cite{Reu}.
\begin{lemma}[\cite{Reu}, Lemma~8.11]\label{lem:reu}
    For $\lambda \vdash n$, let
    \[
    \pi_\lambda(q) = \frac{(q;q)_n}{(1-q^{\lambda_1}) \cdots (1-q^{\lambda_{\ell(\lambda)}})}.
    \]
    If $d \mid n$ and $\omega$ is a primitive $d$-th root of unity, then
    \[
    \pi_\lambda(\omega) = \begin{cases}
        0 & \text{if } \lambda \neq (d^{n/d}), \\
        z_\lambda & \text{if } \lambda = (d^{n/d}).
    \end{cases}
    \]
\end{lemma}
We now proceed with the proof.
\begin{proof}[Proof (of \Cref{thm:super-KW}).]
    Let 
    \[
    f_\lambda(\xb;q,t) = \sum_{\lambda \vdash n+m} s_\lambda(\xb) \sum_{\mathcal{T} \in \SYT_\pm(\lambda)} q^{\maj(\mathcal{T})}t^{\negg(\mathcal{T})}.
    \]
    Then it suffices to show that
    \[
    \ch(\Ltil_{n,m};\xb) = [t^m]\Omega_1^{n+m}f_\lambda(\xb;q,t).
    \]
    By \Cref{thm:s-ps} and \Cref{lem:s-p}, we have
    \begin{align*}
    f_\lambda(\xb;q,t) &= (q;q)_n \sum_{\lambda \vdash n+m} s_\lambda(\xb) \stil_\lambda(1,q,q^2,\ldots;t,tq,tq^2,\ldots) \\
    &= (q;q)_n \sum_{\lambda \vdash n+m} \frac{p_\lambda(\xb)}{z_\lambda} \ptil_\lambda(1,q,q^2,\ldots;t,tq,tq^2,\ldots).
    \end{align*}
    Note that
    \begin{align*}
    \ptil_\lambda(1,q,q^2,\ldots;t,tq,tq^2,\ldots) &= \prod_{j=1}^{\ell(\lambda)} (p_{\lambda_j}(1,q,q^2,\ldots) + (-1)^{\lambda_j+1}p_{\lambda_j}(t,tq,tq^2,\ldots)) \\
    &= \prod_{j=1}^{\ell(\lambda)} \left( \frac{1}{1-q^{\lambda_j}} + (-1)^{\lambda_j+1}\frac{t^{\lambda_j}}{1-q^{\lambda_j}} \right),
    \end{align*}
    so that
    \[
    f_\lambda(\xb;q,t) = \sum_{\lambda \vdash n+m} \frac{p_\lambda(\xb)}{z_\lambda} \pi_\lambda(q) \sum_{J \subseteq [\ell(\lambda)]} \prod_{j \in J} (-1)^{\lambda_j+1}t^{\lambda_j}.
    \]
    Now, for any $(n+m)$-th root of unity $\zeta$, we have by \Cref{lem:reu} that $\pi_\lambda(\zeta) = 0$ unless $\lambda = (d^\frac{n+m}{d})$ and $\zeta$ is a primitive $d$-th root of unity, in which case $\pi_\lambda(\zeta) = z_\lambda$. Thus
    \begin{align*}
    \Omega_1^{n+m}f_\lambda(\xb;q,t) &= \frac{1}{n+m} \sum_{\zeta^{n+m} = 1} \zeta^{-1} f_\lambda(\xb,\zeta,t) \\
    &= \frac{1}{n+m} \sum_{\zeta^{n+m} = 1} \zeta^{-1} \sum_{\lambda \vdash n+m} \frac{p_\lambda(\xb)}{z_\lambda} \pi_\lambda(\zeta) \sum_{J \subseteq [\ell(\lambda)]} \prod_{j \in J} (-1)^{\lambda_j+1}t^{\lambda_j} \\
    &= \frac{1}{n+m} \sum_{d \mid n+m} \sum_{\substack{\zeta \text{ a primitive} \\ \text{$d$-th root of unity}}} \zeta^{-1} p_d(\xb)^\frac{n+m}{d} \sum_{J \subseteq [\frac{n+m}{d}]} (-1)^{d|J| + |J|}t^{d|J|} \\
    &= \frac{1}{n+m} \sum_{d \mid n+m} \mu(d) p_d(\xb)^\frac{n+m}{d} \sum_{J \subseteq [\frac{n+m}{d}]} (-1)^{d|J| + |J|}t^{d|J|}.
    \end{align*}
    When we extract the coefficient of $t^m$, we may restrict the first sum to all $d \mid n+m$ such that also $d \mid m$, so that $d \mid \gcd(n,m)$. Therefore
    \begin{align*}
    [t^m]\Omega_1^{n+m}f_\lambda(q,t) &= \frac{1}{n+m} \sum_{d \mid \gcd(n,m)} \mu(d) p_d(\xb)^\frac{n+m}{d} \sum_{\substack{J \subseteq [\frac{n+m}{d}], \\ |J| = \frac{m}{d}}} (-1)^{m + \frac{m}{d}} \\
    &= \frac{1}{n+m} \sum_{d \mid \gcd(n,m)} \mu(d) p_d(\xb)^\frac{n+m}{d} (-1)^{m + \frac{m}{d}} \binom{\frac{n+m}{d}}{\frac{m}{d}} \\
    &= \ch(\Ltil_{n,m};\xb),
    \end{align*}
    where the last equality follows from \Cref{thm:super-brandt}.
\end{proof}

\section{Further directions}\label{sec:addl}

\subsection{Idempotents}

Klyachko \cite{Kly} introduced a remarkable idempotent
  \[ \mathcal{K}_n = \frac{1}{n} \sum_{\sigma \in S_n} \omega^{\maj(\sigma)} \sigma \]
where $\omega$ is a primitive $n$th root of unity. He showed $V^{\otimes n} \mathcal{K}_n = \Lcal_n(V)$. See \cite[\S4]{Gar} for a combinatorial approach. Separately, Reutenauer noted
  \[ \Trm(V) = S(\Lcal(V)) = \bigoplus_{d=0}^\infty \Lcal(V)^{(d)} \]
where $\Lcal(V)^{(d)}$ is spanned by $d$th powers of elements of $\Lcal(V)$ \cite[\S2]{Reu-idemp}. Reutenauer gave orthogonal idempotents $\rho_n^{(d)} \in \QQ S_n$ projecting $V^{\otimes n}$ onto $\Lcal(V)^{(d)}$ for $1 \leq d \leq n$. See also \cite[Thm.~7.3]{Gar} and the surrounding discussion. It would be interesting to investigate super versions of these idempotents, which we leave as an open problem.

\begin{prob}
    Find super analogues of the Klyachko idempotent $\mathcal{K}_n$ and the Reutenauer idempotents $\rho_n^{(d)}$.
\end{prob}

\subsection{Combinatorial symmetry}

A slight variation on the Kr\'askiewicz--Weyman formula shows that the multiplicity of $S^\lambda$ in $\chi^r\uparrow_{C_n}^{S_n}$ is $\#\{T \in \SYT(\lambda) : \maj(T) \equiv_n r \}$. Since the isomorphism type of $\chi^r\uparrow_{C_n}^{S_n}$ depends only on $n$ and $\gcd(r, n)$, we have the following enumerative corollary. If $\gcd(r, n) = \gcd(s, n)$, then
\begin{equation}\label{eq:sym.1}
  \#\{T \in \SYT(\lambda) : \maj(T) \equiv_n r \} = \#\{T \in \SYT(\lambda) : \maj(T) \equiv_n s \}.
\end{equation}
A combinatorial proof of \eqref{eq:sym.1} is currently unknown. See \cite{AS} for further discussion of this problem.

Similarly, we find the multiplicity of $S^\lambda$ in $(\chi^{\mathrm{cyc}} \otimes \chi^{m/\!/2+r})\uparrow_{C_{n+m}}^{S_{n+m}}$ where
  \[ m/\!/2 \coloneqq \begin{cases}m/2 & \text{if $m$ is even} \\ 0 & \text{if $m$ is odd}\end{cases} \]
is
\begin{equation}
  \#\{\mathcal{T} \in \SYT_\pm(\lambda) : \maj(\mathcal{T}) \equiv_{n+m} r, \negg(\mathcal{T}) = m\}.
\end{equation}
The isomorphism type in this case depends only on $n+m$, $\gcd(r + m/\!/2, n+m)$, and $\binom{n+m}{m}$. In particular, if $\gcd(r + m/\!/2, n+m) = \gcd(s + m/\!/2, n+m)$, then
\begin{equation}\label{eq:sym.2}
\begin{split}
  &\#\{\mathcal{T} \in \SYT_\pm(\lambda) : \maj(\mathcal{T}) \equiv_{n+m} r, \negg(\mathcal{T}) = m\} \\
  =\ &\#\{\mathcal{T} \in \SYT_\pm(\lambda) : \maj(\mathcal{T}) \equiv_{n+m} s, \negg(\mathcal{T}) = m\}.
\end{split}
\end{equation}
and, if $n, m$ are odd,
\begin{equation}\label{eq:sym.3}
\begin{split}
  &\#\{\mathcal{T} \in \SYT_\pm(\lambda) : \maj(\mathcal{T}) \equiv_{n+m} r, \negg(\mathcal{T}) = m\} \\
  =\ &\#\{\mathcal{T} \in \SYT_\pm(\lambda) : \maj(\mathcal{T}) \equiv_{n+m} r, \negg(\mathcal{T}) = n\}.
\end{split}
\end{equation}

\begin{prob}
    Find a combinatorial proof of \eqref{eq:sym.1}, \eqref{eq:sym.2}, or \eqref{eq:sym.3}.
\end{prob}

\subsection{The degree two case}

We conclude by discussing the extension of the only other known case of Thrall's problem to the super setting. Recall that
\[
\Ch(\Lcal_{(2^d)}) = h_d[e_2] = \sum_\mu s_\mu,
\]
where the sum is over all partitions of $2d$ with even column lengths. We have $\Lcal_{(2^d)} = S^d(\Lcal_2)$, so a natural extension of this case to the free Lie superalgebra would consist in determining the Schur expansion of the characters of the super higher Lie modules
\[
\Ltil_A = S^d(\Ltil_{2,0}), \Ltil_B = \bigwedge^d(\Ltil_{1,1}), \text{ and } \Ltil_C = S^d(\Ltil_{0,2}),
\]
where the matrices $A,B$, and $C$ are given by $a_{2,0} = b_{1,1} = c_{0,2} = d$ for any $d \ge 2$ and $a_{i,j} = b_{i,j} = c_{i,j} = 0$ for all other pairs $i,j$. The case of $\Ltil_A$ is classical, and the character of $\Ltil_C$ can be computed similarly by a known plethystic identity, which can also be found in \cite[Ex.~I.8.6(a)]{Mac}:
\[
\Ch(\Ltil_C) = h_d[\Ch(\Ltil_{0,2})] = h_d[h_2] = \sum_{\nu} s_\nu,
\]
where the sum is over all partition $\nu \vdash 2d$ with even parts.

Finally, the character of $\Ltil_B$ can be written as follows:
\[
\Ch(\Ltil_B) = e_d[h_2+e_2] = \sum_{k = 0}^d e_k[h_2]e_{d-k}[e_2].
\]
The Schur expansions of $e_k[h_2]$ and $e_k[e_2]$ for arbitrary $k$ are given in \cite[Ex.~I.8.6(c,d)]{Mac}, so that the Schur expansion of $\Ch(\Ltil_B)$ can be determined via the Littlewood--Richardson rule. However, to our knowledge there is no explicit combinatorial description of the coefficients in this expansion.

